\documentclass[12pt,a4paper]{article}

% --- Packages ---
\usepackage[utf8]{inputenc}
\usepackage[T1]{fontenc}
\usepackage{mathptmx}           % Police Times
\usepackage[french]{babel}
\usepackage{amsmath,amssymb}
\usepackage{graphicx}
\usepackage{booktabs}
\usepackage{multirow}
\usepackage{array}
\usepackage{tabularx}
\usepackage{float}
\usepackage{hyperref}
\usepackage{natbib}
\usepackage[margin=2.5cm]{geometry}
\usepackage{setspace}
\onehalfspacing
\usepackage{xcolor}

% --- Titre ---
\title{Référence Formantique de l'Orchestre :\\
Étude quantitative des convergences spectrales instrumentales\\
et fondements acoustiques des doublures orchestrales}

\author{Yan Maresz\\[6pt]
\small Compositeur, Professeur de nouvelles technologies appliquées à la composition\\
\small Conservatoire National Supérieur de Musique et de Danse de Paris\\
\small \texttt{y.maresz@cnsmdp.fr}
}

\date{}

\begin{document}

\maketitle


\begin{center}
\small
\textit{La page web
\href{https://zseramnay.github.io/Orchestration/Etude-Formants/Versions-html-and-docx/REFERENCE_FORMANTIQUE_COMPLETE.html}{"REFERENCE\_FORMANTIQUE\_COMPLETE"} \\
est un complément indispensable à cet article !}
\end{center}
\vspace{1cm}


% ============================================================
% RÉSUMÉ
% ============================================================
\begin{abstract}
Cette étude présente la première analyse quantitative systématique des formants
spectraux de 30 instruments de l'orchestre symphonique, à partir de
5\,914 échantillons de tenues soutenues issus de la base SOL2020
(IRCAM/Orchidea) et d'une collection complémentaire (Yan\_Adds). Nous
introduisons la mesure du Formant Principal (\textit{Fp}), un centroïde
spectral pondéré en énergie dans une bande fréquentielle limitée, qui capture
l'identité timbrale d'un instrument avec des gains de stabilité allant jusqu'à
10,4$\times$ par rapport aux mesures brutes de pics. En utilisant le
\textit{Fp} et l'analyse des pics formantiques, nous démontrons que les
doublures orchestrales classiques s'expliquent par des convergences
formantiques mesurables. Le résultat central est l'identification d'un cluster
de convergence dans la zone vocalique /o/--/\aa/ (377--506~Hz), regroupant
11~instruments de 4~familles différentes avec des écarts aussi faibles que
$\Delta = 0$~Hz (flûte--hautbois) et $\Delta = 11$~Hz (basson--violon,
cor--alto). Neuf principes d'orchestration acoustique sont formulés et validés
par un taux de concordance multi-sources de 93\,\% (27/29 instruments
comparables). En établissant un pont entre la littérature sur le centroïde
spectral et la fusion timbrale \citep{Sandell1995,Tardieu2012,Goodchild2018}
et la tradition des traités d'orchestration, cette étude fournit un cadre
quantitatif pour l'analyse et la pratique de l'orchestration.

\medskip
\noindent\textbf{Mots-clés :} formants instrumentaux, orchestration, analyse spectrale,
doublures, centroïde spectral, convergence formantique, fusion timbrale, SOL2020, Orchidea
\end{abstract}

% ============================================================
% 1. INTRODUCTION
% ============================================================
\section{Introduction}

Les doublures orchestrales --- l'association de deux instruments ou plus jouant
la même ligne mélodique ou des lignes complémentaires --- constituent l'un des
outils fondamentaux de l'écriture orchestrale depuis le
\textsc{xviii}\textsuperscript{e}~siècle. La littérature des traités
d'orchestration, de Berlioz \citep{Berlioz1844} à Rimsky-Korsakov
\citep{RimskyKorsakov1908}, Koechlin \citep{Koechlin1943}, Piston
\citep{Piston1955} et Adler \citep{Adler1982}, prescrit ces associations sur
la base de l'expérience auditive accumulée, sans quantification acoustique
systématique.

Parallèlement, l'étude psychoacoustique de la fusion timbrale a établi que le
centroïde spectral joue un rôle central dans la détermination de la qualité de
fusion perceptive entre paires d'instruments. Dans son travail fondateur,
Sandell a démontré que la fusion se détériore en fonction de la hauteur du
centroïde composite et de la différence de centroïde entre instruments,
expliquant jusqu'à 63\,\% de la variance des jugements de fusion pour des
dyades séparées par une tierce mineure
\citep{Sandell1991,Sandell1995}. Des recherches ultérieures ont confirmé et
étendu ces résultats : Kendall et Carterette ont montré que la fusion peut être
prédite à partir de la proximité dans l'espace perceptif du timbre
\citep{Kendall1993} ; Tardieu et McAdams ont étudié la perception de la fusion
pour des dyades de sons impulsifs et soutenus \citep{Tardieu2012} ; et Lembke
et McAdams ont démontré que les caractéristiques locales de l'enveloppe
spectrale, incluant les pics formantiques et leurs relations fréquentielles,
prédisent environ 85\,\% de la variance de la fusion timbrale
\citep{Lembke2017}. Plus récemment, Goodchild et McAdams ont proposé une
taxonomie complète des effets de groupement orchestral fondée sur les principes
de l'analyse de scènes auditives \citep{Goodchild2018,McAdams2022}.

Malgré l'ensemble de ces travaux, un fossé subsiste entre la littérature
psychoacoustique sur la fusion timbrale --- qui étudie typiquement un petit
nombre de paires d'instruments en conditions contrôlées --- et la tradition des
traités d'orchestration, qui couvre l'orchestre complet mais manque de
fondements quantitatifs. La présente étude vise à combler ce fossé en
fournissant une analyse systématique des formants spectraux de 30~instruments
d'orchestre, identifiant les convergences acoustiques qui sous-tendent les
doublures prescrites par la littérature des traités.

Les contributions de cet article sont triples : (1)~une base de données
formantiques complète couvrant 30~instruments sur 5\,914~échantillons, validée
contre quatre sources académiques indépendantes ; (2)~l'introduction du Formant
Principal (\textit{Fp}), un centroïde spectral à bande limitée offrant une
caractérisation timbrale considérablement plus stable que l'extraction brute de
pics pour les instruments à large tessiture ; et (3)~l'identification d'un
cluster de convergence inter-familles dans la zone vocalique /o/--/\aa/ qui
explique quantitativement les doublures les plus valorisées de la littérature
orchestrale.

% ============================================================
% 2. ÉTAT DE L'ART
% ============================================================
\section{État de l'art}

\subsection{Descripteurs du timbre et centroïde spectral}

Le centroïde spectral --- le centre de gravité du spectre de puissance --- a
été identifié comme l'un des descripteurs du timbre les plus pertinents
perceptivement depuis les études pionnières d'analyse multidimensionnelle de
Grey \citep{Grey1977}. McAdams et al.\ ont confirmé sa corrélation avec la
dimension de « brillance » des espaces de timbre \citep{McAdams1995}. La
\emph{Timbre Toolbox} \citep{Peeters2011}, développée à l'IRCAM, fournit un
cadre standardisé pour le calcul du centroïde spectral et des descripteurs
associés à partir de signaux musicaux. Dans la norme MPEG-7, le centroïde
spectral est inclus comme descripteur audio de base \citep{MPEG7}.

Une distinction essentielle doit être établie entre le centroïde spectral
\emph{global} (calculé sur l'ensemble du spectre) et les mesures de centroïde
\emph{à bande limitée}. Le centroïde spectral global capture la brillance
d'ensemble mais n'isole pas la structure formantique du corps résonant de
l'instrument. Pour les instruments à large tessiture, où la fréquence
fondamentale balaie une grande plage, le centroïde global suit la hauteur
plutôt que le timbre. La mesure \textit{Fp} introduite dans cette étude
répond à cette limitation en calculant le centroïde dans une bande
fréquentielle spécifique à l'instrument, optimisée pour capturer la résonance
principale (voir section~\ref{sec:fp}).

\subsection{Les formants en acoustique musicale}

Le concept de formants --- des régions de résonance fixes du corps de
l'instrument qui façonnent l'enveloppe spectrale indépendamment de la hauteur
jouée --- trouve son origine dans l'acoustique de la parole \citep{Fant1960} et
a été étendu aux instruments de musique par Backus \citep{Backus1969}, Meyer
\citep{Meyer2009} et Gieseler et al.\ \citep{Gieseler1985}. Les mesures
étendues de Meyer ont établi des correspondances vocaliques pour les timbres
instrumentaux, décrivant le timbre du cor comme « semblable au /o/ » et celui
du hautbois comme « semblable au /a/ ». Cependant, ces caractérisations sont
restées qualitatives, fondées sur des impressions auditives plutôt que sur une
analyse spectrale systématique sur de grands corpus.

Fletcher et Rossing fournissent le fondement en acoustique physique, montrant
comment les propriétés résonantes des corps d'instruments (colonnes d'air,
couplage corde-caisse, pavillon) créent des caractéristiques fixes de
l'enveloppe spectrale \citep{Fletcher1998}. Pour les cordes frottées, le
\emph{Hauptformant} (formant principal) décrit par Meyer correspond à la
résonance principale du corps et est relativement stable sur toute la tessiture
de l'instrument --- une propriété que nous quantifions systématiquement dans
cette étude.

\subsection{La fusion timbrale en orchestration}

La thèse de doctorat de Sandell et sa publication subséquente ont établi le
fondement empirique de la compréhension de la fusion timbrale
\citep{Sandell1991,Sandell1995}. Ses trois expériences avec des dyades
d'instruments aux sonorités naturelles ont montré que : (a)~les instruments à
centroïdes spectraux plus bas fusionnent mieux ; (b)~des différences de
centroïde plus faibles entre instruments améliorent la fusion ; et (c)~la
similarité des caractéristiques d'attaque et la corrélation des enveloppes
d'amplitude contribuent à la fusion. Ces facteurs expliquaient de 51\,\%
(unissons) à 63\,\% (tierces mineures) de la variance des jugements de fusion.

Des travaux plus récents de Lembke et McAdams ont étendu cette approche aux
descriptions d'enveloppe spectrale, démontrant que les caractéristiques au
niveau des formants --- incluant les relations fréquentielles entre maxima
spectraux, leur proéminence et les contraintes du système auditif --- prédisent
environ 85\,\% de la variance de la fusion \citep{Lembke2015,Lembke2017}. Ce
résultat motive l'attention portée dans la présente étude à la convergence
formantique comme prédicteur de la fusion orchestrale.

Le projet ACTOR (\emph{Analysis, Creation, and Teaching of ORchestration}) de
l'Université McGill a produit une taxonomie des effets de groupement orchestral
fondée sur l'analyse de scènes auditives
\citep{Goodchild2018,McAdams2022}, distinguant entre l'augmentation timbrale
(un timbre dominant enrichi par d'autres) et l'émergence timbrale (un nouveau
timbre non identifiable comme l'un de ses constituants). Notre cadre de
convergence formantique fournit des corrélats acoustiques pour ces catégories
perceptives : des formants convergents ($\Delta \leq 80$~Hz) prédisent
l'augmentation/fusion, tandis que des formants divergents ($\Delta > 200$~Hz)
prédisent la complémentarité ou l'émergence.

\subsection{Orchestration computationnelle}

Le système Orchidea \citep{Cella2020}, développé à partir de la base SOL à
l'IRCAM, utilise l'appariement spectral pour trouver des combinaisons
d'instruments approximant un son cible. Notre travail complète cette approche
en fournissant une analyse au niveau des formants qui explique \emph{pourquoi}
certaines combinaisons fonctionnent perceptivement, plutôt que de trouver des
combinaisons par appariement spectral exhaustif.

% ============================================================
% 3. CORPUS ET SOURCES
% ============================================================
\section{Corpus et sources}
\label{sec:corpus}

\subsection{Données primaires}

Le corpus comprend deux sources complémentaires :

\begin{itemize}
\item \textbf{SOL2020 (IRCAM/Orchidea)} --- 12 instruments solistes,
2\,391 échantillons de tenues soutenues. La base SOL2020 fournit des
enveloppes spectrales pré-calculées (\texttt{specenv}, 1\,024 bins) issues
d'enregistrements réalisés dans la chambre anéchoïque de l'IRCAM à
44\,100~Hz / 24 bits. Instruments couverts : piccolo, flûte, hautbois,
clarinette en si$\flat$, basson, cor en fa, trompette en ut, trombone, tuba,
saxophone alto, violon, alto, violoncelle et contrebasse.

\item \textbf{Yan\_Adds} --- 18 instruments supplémentaires (ensembles de
cordes, bois auxiliaires, cuivres graves, variantes avec sourdine),
1\,646 échantillons provenant de la Vienna Symphonic Library (VSL), de Con
Timbre et d'enregistrements personnels. Tous les échantillons ont été
enregistrés en conditions de studio contrôlées à 44\,100~Hz ou 48\,000~Hz
(rééchantillonnés à 44\,100~Hz) et normalisés à $-1$~dBFS crête. Cette
collection étend le corpus aux instruments et configurations non couverts par
SOL2020, incluant la flûte basse, la flûte contrebasse, le cor anglais, la
clarinette en mi$\flat$, la clarinette basse, la clarinette contrebasse, le
trombone basse, le tuba basse, le tuba contrebasse, le contrebasson, et les
sections de cordes (ensemble de violons, ensemble d'altos, ensemble de
violoncelles, ensemble de contrebasses), ainsi que les variantes avec sourdine
pour les cordes, le cor, la trompette, le trombone et le tuba.
\end{itemize}

Au total, 487 lignes de données SOL2020 et 46 lignes Yan\_Adds (fichiers CSV
v3) couvrent les 30 instruments analysés. Toutes les données et scripts
d'analyse sont disponibles dans le dépôt
associé\footnote{\url{https://github.com/zseramnay/Orchestration}}.

\subsection{Sources de validation}

Quatre références académiques indépendantes servent de validation croisée :

\begin{enumerate}
\item \textbf{Backus} \citep{Backus1969} --- mesures ciblées des formants des
cuivres et des bois.
\item \textbf{Gieseler et al.} \citep{Gieseler1985} --- mesures détaillées avec
associations vocaliques explicites pour 13 instruments.
\item \textbf{Meyer} \citep{Meyer2009} --- analyse acoustique avec descriptions
de couleurs vocales et zones vocaliques définies.
\item \textbf{McCarty/CCRMA} \citep{McCarty2003} --- référence directionnelle
avec limitations d'exactitude reconnues par l'auteur.
\end{enumerate}

Le taux de concordance global est de 93\,\% (27 instruments sur 29
comparables). Les deux discordances concernent le piccolo (Backus cite
$\sim$3\,400~Hz, correspondant vraisemblablement au pic de rayonnement plutôt
qu'à F1) et la trompette (forte variabilité de F1 selon les sources due au
comportement spectral dépendant de la dynamique).

% ============================================================
% 4. MÉTHODOLOGIE
% ============================================================
\section{Méthodologie}
\label{sec:methodology}

\subsection{Périmètre de mesure}

Toutes les valeurs formantiques sont mesurées sur des tenues soutenues
(\emph{sustained}) en jeu \emph{ordinario}, excluant les transitoires
d'attaque, les modes de jeu étendus et les dynamiques extrêmes. Techniques
incluses : normal, arco, sustain, ordinario, vibrato, non-vibrato, flautando,
de \emph{pp} à \emph{ff}. Techniques exclues : pizzicato, col legno, tremolo,
trilles, harmoniques artificiels, multiphoniques, flatterzunge, sourdines
(analysées séparément), glissandi, portamento. Ces exclusions garantissent que
les formants extraits reflètent les résonances structurelles de l'instrument
plutôt que des artefacts spectraux transitoires.

\subsection{Pipeline d'extraction formantique}

Le pipeline d'extraction utilise les paramètres suivants :

\begin{itemize}
\item \textbf{Taille FFT :} 4\,096 échantillons à 44\,100~Hz, soit une
résolution fréquentielle de $\sim$10,8~Hz/bin.
\item \textbf{Plage analysée :} 100--3\,000~Hz.
\item \textbf{Données d'entrée :} enveloppes spectrales pré-calculées
(\texttt{specenv}, 1\,024 bins) issues du pipeline Orchidea, qui applique un
fenêtrage de Hann et un lissage en amont.
\item \textbf{Détection de pics :} \texttt{scipy.signal.find\_peaks} avec les
paramètres suivants : \texttt{height}~$= -30$~dB par rapport au maximum
régional, \texttt{distance}~$= 7$ bins ($\approx$70~Hz de distance minimale
inter-pics), \texttt{prominence}~$= 3$~dB.
\end{itemize}

Un choix méthodologique critique est la détection du pic \emph{le plus
proéminent} (amplitude la plus élevée) plutôt que du \emph{premier} pic
(fréquence la plus basse). Cela capture le formant perceptuellement dominant et
évite les résonances basses secondaires qui peuvent être acoustiquement
présentes mais perceptivement subordonnées.

Le pipeline (version v22) a été validé par test aller-retour : les 16 fichiers
CSV instrumentaux ont été régénérés et comparés aux valeurs de référence,
atteignant $\Delta = 0$~Hz sur toutes les entrées.

\subsection{Note terminologique}

Dans cette étude, $F_1, F_2, \ldots F_6$ désignent les pics spectraux
principaux extraits par analyse d'enveloppe spectrale. Ces pics correspondent
aux résonances acoustiques principales du système de l'instrument et sont mis
en regard des voyelles cardinales par analogie perceptive, sans prétendre à une
équivalence stricte avec les formants du tractus vocal humain au sens de Fant
\citep{Fant1960}. Pour les instruments à tuyau cylindrique (clarinette, flûte)
où aucun formant fixe n'existe au sens strict, $F_1$ désigne la zone d'énergie
spectrale dominante.

\subsection{Le Formant Principal (\textit{Fp})}
\label{sec:fp}

Pour les instruments à large tessiture, le pic $F_1$ est très instable d'un
registre à l'autre. Le violon, par exemple, présente un $F_1$ avec
$\sigma = 651$~Hz --- soit une variation dépassant une octave. Cette
instabilité rend $F_1$ inexploitable comme signature timbrale.

La mesure \textit{Fp} (Formant Principal) résout ce problème. Elle est définie
comme le centroïde spectral pondéré en énergie dans une bande fréquentielle
spécifique à l'instrument :

\begin{equation}
F_p = \frac{\sum_{i} f_i \cdot A_i}{\sum_{i} A_i}
\label{eq:fp}
\end{equation}

\noindent où $f_i$ est la fréquence du bin $i$, $A_i$ son amplitude (échelle
linéaire), et la somme porte sur une bande choisie par instrument.

\paragraph{Relation avec le centroïde spectral global.}
Le \textit{Fp} diffère du centroïde spectral standard \citep{Peeters2011} sur
deux points : (1)~il est calculé dans une bande fréquentielle restreinte plutôt
que sur l'ensemble du spectre, et (2)~les limites de la bande sont optimisées
par instrument pour capturer la résonance principale du corps
(\emph{Hauptformant} dans la terminologie de Meyer \citep{Meyer2009}). Le
centroïde spectral global, appliqué à l'ensemble du spectre, suit la brillance
d'ensemble et est fortement influencé par la fréquence fondamentale et par la
distribution d'énergie dans les hautes fréquences. À l'inverse, le \textit{Fp}
isole la signature résonante du corps de l'instrument, qui est la
caractéristique acoustiquement pertinente pour prédire la fusion timbrale dans
les doublures.

\paragraph{Sélection de la bande.}
Les bandes fréquentielles ont été déterminées par instrument en identifiant la
région d'énergie maximale de l'enveloppe spectrale sur l'ensemble des registres
en jeu \emph{ordinario}. Les bandes typiques sont 600--1\,400~Hz pour les
cordes et le cor, et 1\,000--2\,000~Hz pour la flûte, le hautbois et la
clarinette. Les limites de bande ont été fixées de manière à englober les
points à $-6$~dB du pic principal de l'enveloppe spectrale moyennée sur tous
les échantillons d'un instrument donné.

\paragraph{Gain de stabilité.}
Le tableau~\ref{tab:fp_stability} montre l'amélioration de stabilité du
\textit{Fp} par rapport à $F_2$ pour quatre instruments représentatifs. Le gain
va de 4,7$\times$ (clarinette) à 10,4$\times$ (trompette). Vingt-deux des
30 instruments du corpus bénéficient d'une mesure \textit{Fp} validée.

\begin{table}[htbp]
\centering
\caption{Comparaison de stabilité : \textit{Fp} vs.\ $F_2$ (ordinario).}
\label{tab:fp_stability}
\begin{tabular}{lrrrrc}
\toprule
Instrument & $F_p$ (Hz) & $\sigma_{F_p}$ & $F_2$ (Hz) & $\sigma_{F_2}$ & Gain \\
\midrule
Violon     & 893  & 92  & 1\,518 & 651   & 7,1$\times$ \\
Trompette  & 1\,046 & 98  & 1\,324 & 1\,018 & 10,4$\times$ \\
Clarinette si$\flat$ & 1\,412 & 169 & 1\,130 & 795 & 4,7$\times$ \\
Cor        & 738  & 112 & 1\,106 & 734   & 6,5$\times$ \\
\bottomrule
\end{tabular}
\end{table}

\paragraph{$\sigma(F_p)$ comme indicateur de sélectivité spectrale.}
L'écart-type $\sigma(F_p)$ reflète indirectement la largeur de bande de la
résonance principale : un $\sigma$ faible indique un formant étroit et bien
défini (facteur $Q$ élevé), tandis qu'un $\sigma$ élevé indique un formant
large et diffus. Le violon ($\sigma = 92$~Hz) possède ainsi un timbre
spectralement plus focalisé que la clarinette si$\flat$ ($\sigma = 169$~Hz),
ce qui est cohérent avec l'expérience perceptive.

\subsection{Deux représentations complémentaires par registre}

Pour chaque instrument principal, le document de référence fournit deux
représentations complémentaires calculées registre par registre (les registres
étant définis selon les nomenclatures instrumentales traditionnelles) :

\begin{enumerate}
\item \textbf{Profil formantique moyen :} extraction des formants $F_1$--$F_6$
sur chaque échantillon individuel, puis calcul de la médiane par registre. Ce
profil donne les positions de référence de chaque pic spectral.

\item \textbf{Carte spectrale vocalique :} moyenne de toutes les enveloppes
spectrales du registre, affichée sur un axe fréquentiel logarithmique avec
superposition des zones vocaliques (/u/, /o/, /a/, /e/) et enveloppe cepstrale
normalisée. Cette carte montre la forme réelle du spectre moyen avec ses zones
de résonance, ses vallées et le contexte vocalique complet.
\end{enumerate}

Les deux approches convergent pour les formants bien définis (instruments à
anche, cuivres) mais peuvent diverger pour les instruments à grande tessiture
(cordes, flûtes, clarinettes) où la moyenne de l'enveloppe lisse les pics
individuels. Le profil formantique donne les positions de référence ; la carte
spectrale donne le contexte acoustique complet.

Cette double analyse par registre résout le problème de la variabilité
spectrale des instruments à large tessiture : plutôt que de se limiter à une
valeur $F_1$ globale peu informative, l'examen registre par registre révèle des
convergences formantiques qui n'apparaissent que lorsque les instruments jouent
dans un registre spécifique (voir section~\ref{sec:principles}, principe~8).

% ============================================================
% 5. CORRESPONDANCE VOYELLES–FRÉQUENCES
% ============================================================
\section{Correspondance voyelles--fréquences}

Le système vocalique fournit un cadre intuitif pour décrire le timbre
instrumental. D'après Meyer \citep{Meyer2009} et Gieseler et al.\
\citep{Gieseler1985}, les correspondances sont présentées dans le
tableau~\ref{tab:vowels}.

\begin{table}[htbp]
\centering
\caption{Zones vocaliques et instruments représentatifs ($F_1$ ou $F_p$).}
\label{tab:vowels}
\begin{tabular}{llll}
\toprule
Voyelle & Fréquence (Hz) & Perception & Instruments représentatifs \\
\midrule
/u/ & 200--400 & Profondeur  & Cb (172), Tuba (226), Cbn (226), Trb (237) \\
/o/ & 400--600 & Plénitude   & Alto (377), Cor (388), Sax a. (398), C.a. (452), \\
    &          &             & Cl (463), Bn (495), Vln (506) \\
/\aa/ & 600--800 & Chaleur   & Cl$^{Mib}$ (678), Fl (743), Htb (743), Tpt (786) \\
/a/ & 800--1\,250 & Puissance & Cor$_{Fp}$ (738), Vln$_{Fp}$ (893), Tpt$_{Fp}$ (1\,046) \\
/e/ & 1\,250--2\,600 & Clarté & Cl$_{Fp}$ (1\,412), Fl$_{Fp}$ (1\,535), Htb$_{Fp}$ (1\,485) \\
/i/ & $>$2\,600 & Brillance & Picc (1\,992 $F_2$), Tpt+harmon (2\,358) \\
\bottomrule
\end{tabular}
\end{table}

% ============================================================
% 6. RÉSULTATS PAR FAMILLE INSTRUMENTALE
% ============================================================
\section{Résultats par famille instrumentale}

\subsection{Les bois (plage : 226--2\,638~Hz)}

La famille des bois présente la plus grande diversité timbrale de l'orchestre.
Les anches doubles (hautbois, cor anglais, basson, contrebasson) partagent une
zone formantique commune autour de 900--1\,200~Hz avec une coloration /a/
caractéristique. Les clarinettes (tuyau cylindrique, harmoniques impairs) n'ont
pas de formant fixe au sens strict : le pic spectral suit la fondamentale ou le
3\textsuperscript{e} harmonique. Les flûtes (tuyau ouvert) présentent un
spectre décroissant de façon approximativement linéaire au-dessus de la
fondamentale.

Découvertes clés : le cor anglais ($F_1 = 452$~Hz) et le basson
($F_1 = 495$~Hz) se situent dans la zone /o/, aux côtés du cor (388~Hz) et
du violon (506~Hz). Le cor anglais et la clarinette si$\flat$ convergent à
$\Delta = 11$~Hz (452 vs.\ 463~Hz). Le basson et le violon convergent à
$\Delta = 11$~Hz (495 vs.\ 506~Hz).

\subsection{Les cuivres (plage : 162--2\,358~Hz)}

Les cuivres présentent des formants bien définis et une grande stabilité
spectrale inter-dynamiques (sauf la trompette). Le cor en fa
($F_1 = 388$~Hz, zone /o/) est l'instrument cuivre le plus polyvalent pour
les doublures, avec un accord unanime des quatre sources académiques. La
trompette en ut présente un $F_1$ extrêmement variable ($\sigma = 642$~Hz),
mais son \textit{Fp} à 1\,046~Hz est 10,4$\times$ plus stable que $F_2$.
Dans le grave, le trombone (237~Hz), le tuba basse (226~Hz) et le tuba
contrebasse (226~Hz) forment un cluster /u/ très serré
($\Delta \leq 11$~Hz), rejoints par le contrebasson (226~Hz également).

\subsection{Les saxophones}

Le saxophone alto ($F_1 = 398$~Hz) s'inscrit pleinement dans la zone de
convergence /o/--/\aa/, aux côtés du cor ($\Delta = 10$~Hz) et de l'alto
($\Delta = 21$~Hz). C'est le seul saxophone présent dans le corpus SOL2020 ;
les saxophones soprano, ténor et baryton constituent une extension future
prioritaire.

\subsection{Les cordes (plage : 172--1\,518~Hz)}

La famille des cordes présente la plus grande variabilité spectrale de
l'orchestre, compensée par la mesure \textit{Fp}. Le violon
($F_1 = 506$~Hz, /o/) et le basson ($F_1 = 495$~Hz) convergent à
$\Delta = 11$~Hz. Le violoncelle rejoint cette convergence via son
$F_2 = 506$~Hz ($\approx F_1$ du basson).

L'effet de section est clairement quantifié : l'ensemble ne déplace pas $F_1$
mais lisse les harmoniques supérieurs ($F_2$ du violon solo : 1\,518~Hz
$\rightarrow$ ensemble : 1\,163~Hz, soit $-23$\,\%). Le résultat perceptif est
un timbre plus fondu --- c'est la différence de texture entre un quatuor et une
section orchestrale. Cet effet de lissage est cohérent avec la constatation de
Sandell selon laquelle des centroïdes composites plus bas améliorent la fusion
\citep{Sandell1995}.

% ============================================================
% 7. CLUSTER DE CONVERGENCE
% ============================================================
\section{Le cluster de convergence /o/--/\aa/}

Le résultat central de cette étude est l'identification d'un cluster de
convergence dans la zone vocalique /o/--/\aa/ (377--506~Hz), regroupant
11 instruments de 4 familles différentes dans un espace de seulement 129~Hz
(tableau~\ref{tab:cluster}).

\begin{table}[htbp]
\centering
\caption{Les 6 instruments les plus proches du centre du cluster.}
\label{tab:cluster}
\begin{tabular}{llrr}
\toprule
Instrument & Famille & $F_1$ (Hz) & $\Delta$ / Violon \\
\midrule
Alto         & Cordes      & 377 & 129~Hz \\
Cor          & Cuivres     & 388 & 118~Hz \\
Sax alto     & Saxophones  & 398 & 108~Hz \\
Cor anglais  & Bois        & 452 & 54~Hz \\
Clarinette si$\flat$ & Bois & 463 & 43~Hz \\
Basson       & Bois        & 495 & 11~Hz \\
Violon       & Cordes      & 506 & --- \\
\bottomrule
\end{tabular}
\end{table}

Les doublures les plus serrées --- basson\,+\,violon ($\Delta = 11$~Hz),
cor\,+\,alto ($\Delta = 11$~Hz) --- correspondent précisément aux associations
les plus valorisées de la littérature orchestrale. En termes de distance sur
l'échelle de Bark \citep{Zwicker1961}, $\Delta = 11$~Hz à 500~Hz correspond à
environ 0,1~Bark, bien en deçà de la bande critique pour l'intégration
spectrale.

\subsection{Doublures vérifiées}

Le tableau~\ref{tab:doublings} présente les doublures principales classées par
$\Delta$ croissant.

\begin{table}[htbp]
\centering
\caption{Principales doublures vérifiées (sélection).}
\label{tab:doublings}
\begin{tabular}{llrrrl}
\toprule
Instr.~A & Instr.~B & $F_1$ A & $F_1$ B & $\Delta$ & Qualité \\
\midrule
Flûte & Hautbois & 743 & 743 & 0 & Quasi-parfaite $\star$ \\
Tuba CB & Contrebasson & 226 & 226 & 0 & Quasi-parfaite $\star$ \\
Cor & Alto & 388 & 377 & 11 & Quasi-parfaite \\
Cor angl. & Cl si$\flat$ & 452 & 463 & 11 & Quasi-parfaite \\
Violon & Basson & 506 & 495 & 11 & Quasi-parfaite \\
Trombone & Tuba & 237 & 226 & 11 & Quasi-parfaite \\
Cor angl. & Cor & 452 & 388 & 64 & Excellente \\
Trompette & Hautbois & 786 & 743 & 43 & Excellente \\
Hautbois & Basson & 743 & 495 & 248 & Complémentaire \\
Trompette & Trombone & 786 & 237 & 549 & Complémentaire \\
\bottomrule
\end{tabular}
\end{table}

L'analyse par registre révèle en outre des convergences invisibles dans les
données globales. Par exemple, la clarinette si$\flat$ au chalumeau (ré3--ré4,
$F_1 = 215$~Hz) converge exactement ($\Delta = 0$~Hz) avec le trombone au
registre grave ; le cor au médium ($F_1 = 323$~Hz) converge avec l'alto au
médium --- confirmant acoustiquement la doublure emblématique de Brahms.

Il est essentiel de souligner que ces convergences sont valables aussi bien à
l'unisson qu'à l'octave ou à tout autre intervalle. La convergence formantique
est une propriété du timbre et non de la hauteur : deux instruments dont les
formants coïncident fusionnent quelle que soit la distance en hauteur entre les
notes jouées. C'est pourquoi les doublures violon--basson et cor--alto
fonctionnent à travers les registres sans perdre leur efficacité acoustique.

% ============================================================
% 8. NEUF PRINCIPES
% ============================================================
\section{Neuf principes d'orchestration acoustique}
\label{sec:principles}

L'analyse des données conduit à formuler neuf principes fondamentaux.

\paragraph{1.\ Convergence formantique (fusion).}
Quand deux instruments partagent un $F_1$ similaire ($\Delta \leq 80$~Hz),
leurs timbres fusionnent : l'oreille perçoit une couleur homogène plutôt que
deux voix distinctes. Seuils mesurés : $\Delta = 0$ (unisson formantique),
$\Delta \leq 30$~Hz (fusion quasi-parfaite, $\leq 0{,}3$~Bark),
$\Delta$ 30--80~Hz (fusion efficace, 0,3--0,7~Bark). Ce résultat est cohérent
avec l'observation de Sandell selon laquelle des différences de centroïde plus
faibles améliorent la fusion \citep{Sandell1995} et avec la théorie des bandes
critiques \citep{Zwicker1961}.

\paragraph{2.\ Complémentarité spectrale (enrichissement).}
Quand $\Delta > 200$~Hz ($> 1{,}5$~Bark), les spectres se complètent sans se
masquer mutuellement, produisant un enrichissement large bande inaccessible à
chaque instrument seul. Exemple : trompette (786~Hz) + trombone (237~Hz),
$\Delta = 549$~Hz.

\paragraph{3.\ Effet de section (homogénéisation).}
Le passage du soliste à l'ensemble ne déplace pas $F_1$ mais lisse les
harmoniques supérieurs : $F_2$ du violon baisse de 23\,\% en section, $F_3$ de
16\,\%. Le résultat perceptif est un timbre plus fondu. Ce résultat est
cohérent avec la prédiction que des centroïdes composites plus bas améliorent
la fusion \citep{Sandell1995}.

\paragraph{4.\ La sourdine comme transposition timbrale.}
Pour les cordes et le cor, la sourdine abaisse systématiquement $F_1$ de 6 à
40\,\%. Pour les cuivres, l'effet dépend du type de sourdine : les sourdines
cup et harmon produisent l'effet opposé --- la sourdine harmon de la trompette
déplace $F_1$ de la zone /a/ à la zone /i/ (+200\,\%), le cas le plus
spectaculaire du corpus.

\paragraph{5.\ Familles formantiques transversales.}
Les regroupements formantiques traversent les frontières organologiques : la
famille /u/ (contrebasse, tuba, contrebasson, trombone), la famille /o/--/\aa/
(alto, cor, saxophone alto, cor anglais, clarinette, basson, violon) et la
famille /a/ (via les valeurs \textit{Fp}) sont plus pertinentes que les
familles instrumentales traditionnelles pour prédire les fusions timbrales.

\paragraph{6.\ La zone /o/--/\aa/ comme attracteur de convergence.}
La concentration de 11 instruments dans 129~Hz autour de 377--506~Hz correspond
au centre géométrique du spectre perceptif humain (voix parlée, bande de
présence). Cette convergence est cohérente avec l'hypothèse que la palette
orchestrale s'est historiquement construite autour de la région spectrale de
sensibilité perceptive maximale, bien que l'établissement d'une relation
causale nécessiterait des investigations historiques et perceptives
complémentaires.

\paragraph{7.\ Rôle du \textit{Fp} dans les doublures à large tessiture.}
Pour les instruments à grande tessiture (violon, trompette, clarinette), le
\textit{Fp} remplace avantageusement $F_1$ ou $F_2$ dans l'analyse des
doublures. Le gain de stabilité atteint 10,4$\times$ pour la trompette.

\paragraph{8.\ Convergences par registre.}
L'analyse par registre met en lumière des convergences invisibles dans les
données globales : $\Delta = 0$~Hz entre clarinette si$\flat$ (chalumeau) et
trombone (grave) ; entre cor (médium) et alto (médium) ; entre cor (médium) et
trompette (médium).

\paragraph{9.\ Hiérarchie de la fusion.}
Les données établissent une hiérarchie quantitative :
$\Delta \leq 30$~Hz (fusion quasi-parfaite, $\leq 0{,}3$~Bark),
$\Delta$ 30--80~Hz (fusion efficace, 0,3--0,7~Bark),
$\Delta$ 80--200~Hz (proximité, 0,7--1,5~Bark),
$\Delta \geq 200$~Hz (complémentarité spectrale, $\geq 1{,}5$~Bark).

% ============================================================
% 9. DISCUSSION
% ============================================================
\section{Discussion}

\subsection{Validité du \textit{Fp}}

Le \textit{Fp}, centroïde spectral à bande limitée, s'avère être la mesure la
plus pertinente pour caractériser le timbre global d'un instrument à large
tessiture. Il surpasse considérablement $F_1$ et $F_2$ en stabilité tout en
correspondant au \emph{Hauptformant} décrit par Meyer pour les cordes
\citep{Meyer2009}.

Alors que le centroïde spectral global \citep{Peeters2011} capture la
brillance, le \textit{Fp} isole l'identité résonante de l'instrument --- la
région fréquentielle où le corps (ou la colonne d'air) impose sa signature
timbrale indépendamment de la hauteur. Cette distinction est analogue à la
différence entre la « fréquence fondamentale de la voix parlée » et les
« fréquences formantiques » en phonétique : toutes deux sont des centroïdes
spectraux, mais mesurés sur des plages fréquentielles différentes et capturant
des attributs perceptifs distincts.

\subsection{Liens avec la recherche sur la fusion timbrale}

Les convergences formantiques identifiées dans cette étude fournissent des
corrélats acoustiques directs pour les phénomènes de fusion étudiés par Sandell
\citep{Sandell1995}, Kendall et Carterette \citep{Kendall1993}, et Lembke et
McAdams \citep{Lembke2017}. En particulier :

\begin{itemize}
\item Le seuil de $\Delta \leq 80$~Hz pour la fusion correspond à une distance
sur l'échelle de Bark de $\leq 0{,}7$~Bark à 500~Hz, cohérent avec les
contraintes de bande critique sur l'intégration spectrale \citep{Zwicker1961}.

\item Le constat que les instruments à formants bas (cor, alto, basson dans la
zone /o/) participent aux doublures les plus efficaces est cohérent avec la
découverte de Sandell selon laquelle la fusion s'améliore avec des centroïdes
composites plus bas.

\item L'effet de section (lissage des harmoniques supérieurs sans déplacement
de $F_1$) fournit un mécanisme pour l'observation perceptive bien connue selon
laquelle les sections orchestrales fusionnent mieux que les instruments
solistes.
\end{itemize}

\subsection{Limites}

Plusieurs limites doivent être reconnues :

\begin{enumerate}
\item \textbf{Couverture du corpus :} seul un saxophone (alto) est inclus. Les
saxophones soprano, ténor et baryton constituent une extension prioritaire.

\item \textbf{Sources hétérogènes :} les échantillons Yan\_Adds proviennent de
conditions d'enregistrement diverses. Bien que tous aient été enregistrés en
conditions de studio ou anéchoïques, des variations de placement de microphone
et de caractéristiques de salle pourraient introduire des biais systématiques.
Le taux de concordance de 93\,\% avec des sources académiques indépendantes
fournit une validation indirecte.

\item \textbf{Absence de validation perceptive :} le lien entre convergence
formantique et fusion perçue est soutenu par la littérature psychoacoustique
existante \citep{Sandell1995,Lembke2017} mais n'a pas été testé
expérimentalement avec les paires d'instruments et les valeurs de $\Delta$
spécifiques rapportées ici. Une expérience d'écoute mesurant les jugements de
fusion en fonction de la distance $F_1$ pour les doublures du
tableau~\ref{tab:doublings} serait une prochaine étape précieuse.

\item \textbf{Tenues soutenues uniquement :} les transitoires d'attaque, les
modes de jeu étendus et les dynamiques extrêmes, qui modifient radicalement le
spectre, ne sont pas couverts. Des recherches antérieures ont montré que la
similarité d'attaque est un prédicteur significatif de la fusion
\citep{Sandell1995}, suggérant qu'un modèle complet de la fusion orchestrale
devrait incorporer des caractéristiques temporelles aussi bien que spectrales.

\item \textbf{Sélection de bande pour le \textit{Fp} :} les bandes
fréquentielles spécifiques par instrument ont été déterminées par analyse
d'enveloppe spectrale et inspection visuelle, introduisant un degré de
subjectivité. Une procédure d'optimisation automatisée (par exemple,
minimisation de la variance inter-registres du centroïde) pourrait améliorer la
reproductibilité.

\item \textbf{Tests statistiques :} le taux de concordance (93\,\%) est un
ratio simple. Les travaux futurs devraient inclure des intervalles de confiance
par bootstrap pour les estimations de $F_1$ et \textit{Fp} et des tests de
permutation pour la significativité du cluster /o/--/\aa/ par rapport à un
modèle nul de formants distribués aléatoirement.
\end{enumerate}

\subsection{Implications pour la composition}

L'identification quantitative des familles formantiques transversales offre aux
compositeurs un outil nouveau : la possibilité de prédire les fusions et
complémentarités timbrales avant même d'écrire, en choisissant les instruments
selon leur zone vocalique plutôt que selon leur famille organologique. Le seuil
de $\Delta \leq 80$~Hz pour la fusion et de $\Delta \geq 200$~Hz pour
l'enrichissement spectral fournit des critères opérationnels qui complètent,
plutôt qu'ils ne remplacent, l'oreille du compositeur.

% ============================================================
% 10. CONCLUSION
% ============================================================
\section{Conclusion}

Cette étude établit que les doublures orchestrales classiques reposent sur des
convergences formantiques mesurables. Le cluster de convergence /o/--/\aa/
(377--506~Hz), regroupant 11 instruments de 4 familles, constitue le noyau
acoustique de l'orchestre symphonique. L'introduction du \textit{Fp} (centroïde
spectral à bande limitée) comme métrique timbrale, combinée à l'analyse par
registre et à la validation multi-sources (93\,\% de concordance), fournit un
cadre quantitatif reliant la tradition des traités d'orchestration à la
littérature psychoacoustique sur la fusion timbrale.

Le document complet de référence, incluant les profils formantiques détaillés,
les cartes spectrales vocaliques par registre et les matrices de convergence
pour les 30 instruments, est disponible en ligne à l'adresse
\url{https://zseramnay.github.io/Orchestration/}.

\subsection*{Déclaration de disponibilité des données}
Toutes les données (fichiers CSV), scripts d'analyse (Python) et le document de
référence complet sont publiquement accessibles à l'adresse
\url{https://github.com/zseramnay/Orchestration}.

\subsection*{Conflit d'intérêts}
L'auteur déclare n'avoir aucun conflit d'intérêts.

\subsection*{Remerciements}
L'auteur remercie l'équipe Orchidea de l'IRCAM pour la mise à disposition de la
base de données SOL2020.

% ============================================================
% RÉFÉRENCES
% ============================================================
\bibliographystyle{apalike}

\begin{thebibliography}{99}

\bibitem[Adler, 1982]{Adler1982}
Adler, S. (1982).
\newblock \emph{The Study of Orchestration}.
\newblock W.W. Norton \& Company, New York.

\bibitem[Backus, 1969]{Backus1969}
Backus, J. (1969).
\newblock \emph{The Acoustical Foundations of Music}.
\newblock W.W. Norton \& Company, New York.

\bibitem[Berlioz, 1844]{Berlioz1844}
Berlioz, H. (1844/1855).
\newblock \emph{Grand traité d'instrumentation et d'orchestration modernes}.
\newblock Schonenberger, Paris.

\bibitem[Cella et~al., 2020]{Cella2020}
Cella, C.-E., Musicé, E. et Musicé, C. (2020).
\newblock OrchideaSOL: a dataset of extended instrumental techniques for
  computer-aided orchestration.
\newblock In \emph{Proceedings of the International Computer Music Conference}.

\bibitem[Fant, 1960]{Fant1960}
Fant, G. (1960).
\newblock \emph{Acoustic Theory of Speech Production}.
\newblock Mouton, La Haye.

\bibitem[Fletcher et Rossing, 1998]{Fletcher1998}
Fletcher, N.~H. et Rossing, T.~D. (1998).
\newblock \emph{The Physics of Musical Instruments}.
\newblock Springer-Verlag, 2\textsuperscript{e} édition.

\bibitem[Gieseler et~al., 1985]{Gieseler1985}
Gieseler, W., Lombardi, L. et Weyer, R.-D. (1985).
\newblock \emph{Instrumentation in der Musik des 20.\ Jahrhunderts}.
\newblock Breitkopf \& H{\"a}rtel, Wiesbaden.

\bibitem[Goodchild et McAdams, 2018]{Goodchild2018}
Goodchild, M. et McAdams, S. (2018).
\newblock Perceptual processes in orchestration.
\newblock In Siedenburg, K., Saitis, C., McAdams, S., Popper, A.~N. et
  Fay, R.~R., éditeurs, \emph{Timbre: Acoustics, Perception, and Cognition},
  pages 247--271. Springer.

\bibitem[Grey, 1977]{Grey1977}
Grey, J.~M. (1977).
\newblock Multidimensional perceptual scaling of musical timbres.
\newblock \emph{Journal of the Acoustical Society of America}, 61(5):1270--1277.

\bibitem[{ISO/IEC}, 2002]{MPEG7}
{ISO/IEC} (2002).
\newblock {ISO/IEC 15938-4 : Information technology --- Multimedia content
  description interface --- Part 4: Audio}.

\bibitem[Kendall et Carterette, 1993]{Kendall1993}
Kendall, R.~A. et Carterette, E.~C. (1993).
\newblock Identification and blend of timbres as a basis for orchestration.
\newblock \emph{Contemporary Music Review}, 9(1--2):51--67.

\bibitem[Koechlin, 1943]{Koechlin1943}
Koechlin, C. (1943).
\newblock \emph{Traité de l'orchestration} (4 volumes).
\newblock Max Eschig, Paris.

\bibitem[Lembke et McAdams, 2015]{Lembke2015}
Lembke, S.-A. et McAdams, S. (2015).
\newblock Spectral-envelope characteristics as predictors of blend in
  orchestral instrument pairs.
\newblock In \emph{Proceedings of the International Symposium on Musical
  Acoustics (ISMA)}.

\bibitem[Lembke et~al., 2017]{Lembke2017}
Lembke, S.-A., Parker, K., Narmour, E. et McAdams, S. (2017).
\newblock Predicting blend between orchestral timbres using generalized
  spectral-envelope descriptions.
\newblock \emph{Journal of New Music Research}, 46(4):1--18.

\bibitem[McAdams et~al., 1995]{McAdams1995}
McAdams, S., Winsberg, S., Donnadieu, S., De~Soete, G. et Krimphoff, J.
  (1995).
\newblock Perceptual scaling of synthesized musical timbres: Common dimensions,
  specificities, and latent subject classes.
\newblock \emph{Psychological Research}, 58(3):177--192.

\bibitem[McAdams, 2022]{McAdams2022}
McAdams, S. (2022).
\newblock A taxonomy of orchestral grouping effects derived from principles of
  auditory perception.
\newblock \emph{Music Theory Online}, 28(3).

\bibitem[McCarty, 2003]{McCarty2003}
McCarty, M. (2003).
\newblock Vowel space chart for orchestral instruments.
\newblock CCRMA, Stanford University.

\bibitem[Meyer, 2009]{Meyer2009}
Meyer, J. (2009).
\newblock \emph{Acoustics and the Performance of Music}.
\newblock Springer-Verlag, 5\textsuperscript{e} édition.

\bibitem[Peeters et~al., 2011]{Peeters2011}
Peeters, G., Giordano, B.~L., Susini, P., Misdariis, N. et McAdams, S.
  (2011).
\newblock The {T}imbre {T}oolbox: Extracting audio descriptors from musical
  signals.
\newblock \emph{Journal of the Acoustical Society of America},
  130(5):2902--2916.

\bibitem[Piston, 1955]{Piston1955}
Piston, W. (1955).
\newblock \emph{Orchestration}.
\newblock W.W. Norton \& Company, New York.

\bibitem[Rimsky-Korsakov, 1908]{RimskyKorsakov1908}
Rimsky-Korsakov, N. (1908).
\newblock \emph{Principes d'orchestration} (posthume).
\newblock Édition russe de musique.

\bibitem[Sandell, 1991]{Sandell1991}
Sandell, G.~J. (1991).
\newblock \emph{Concurrent Timbres in Orchestration: A Perceptual Study of
  Factors Determining ``Blend''}.
\newblock Thèse de doctorat, Northwestern University.

\bibitem[Sandell, 1995]{Sandell1995}
Sandell, G.~J. (1995).
\newblock Roles for spectral centroid and other factors in determining
  ``blended'' instrument pairings in orchestration.
\newblock \emph{Music Perception}, 13(2):209--246.

\bibitem[Tardieu et McAdams, 2012]{Tardieu2012}
Tardieu, D. et McAdams, S. (2012).
\newblock Perception of dyads of impulsive and sustained instrument sounds.
\newblock \emph{Music Perception}, 30(2):117--128.

\bibitem[Zwicker et Fastl, 1999]{Zwicker1961}
Zwicker, E. et Fastl, H. (1999).
\newblock \emph{Psychoacoustics: Facts and Models}.
\newblock Springer-Verlag, 3\textsuperscript{e} édition.

\end{thebibliography}

\end{document}
